\paragraph{Scale Connection, Holonomy, and Time Characteristic Classes}

\begin{theorem}[Holonomy Criterion for the Scale--Time Connection: Global Unified Time If and Only If the Discrete Curvature Group Vanishes]\label{thm:xi-scale-time-connection-holonomy-criterion}
Let the resolution tower \((X_m,p_{m+1,m})\) induce a groupoid \(\mathcal G_m\) at each level, with a given strictly additive time cocycle \(\ell_m:\mathcal G_m\to\mathbb Z\).
Define the scale lifting defect (the \(1\)-form of the discrete connection)
\[
a_m(g):=\ell_{m+1}(g)-\ell_m\!\bigl(p_{m+1,m}(g)\bigr).
\]
For any commutable "scale--time square" \(g\Rightarrow g'\) (i.e., a pair of parallel morphisms confluent in the certificate sense), define its holonomy
\[
\mathrm{Hol}_m(g\Rightarrow g'):=a_m(g)-a_m(g')\in\mathbb Z.
\]
Let \(\mathfrak H\le \mathbb Z\) be the subgroup generated by all \(\mathrm{Hol}_m(g\Rightarrow g')\).
Then a global unified time exists (equivalently, there exists a family of gauge corrections making the time cocycles at each level strictly compatible under projection) if and only if \(\mathfrak H=\{0\}\).
Moreover, when \(\mathfrak H\neq 0\), the central rank of the minimal central extension required for unified time equals \(\mathrm{rank}(\mathfrak H)\).
\end{theorem}

\begin{proof}
In the bicomplex of Theorem \ref{thm:xi-scale-time-bicomplex-obstruction}, \(a_m=\partial\ell\) is the vertical differential, and \(\mathrm{Hol}_m\) is precisely the value of \(\delta a_m=\delta\partial\ell\) on fundamental commutative squares.
Therefore \(\mathfrak H\) is exactly the observable generating subgroup of the obstruction class \(\mathfrak o\in E_2^{2,0}\); \(\mathfrak H=\{0\}\) is equivalent to \(\mathfrak o=0\), which in turn is equivalent to the existence of a global unified time.
The minimal central rank conclusion follows similarly from the rank of the abelian subgroup generated by \(\mathfrak o\) (Theorem \ref{thm:xi-scale-time-bicomplex-obstruction}).
\end{proof}

\begin{theorem}[Integrality of the Time Characteristic Class: Composition Anomaly and Mixed Curvature Unified into a Single Integer \(H^2\) Class]\label{thm:xi-time-characteristic-class-integrality}
Let \(\mathcal B\) be a \(2\)-category carrying two types of composable structures (e.g., an additive layer \(\circ_+\) and a multiplicative layer \(\circ_\times\)), and assume that each layer is equipped with an integer-valued length cocycle such that:
the composition anomaly is represented by the \(2\)-cocycle \(\alpha\) (as in Proposition \ref{prop:xi-time-length-cocycle}), and the alternating square defect is represented by the mixed curvature \(2\)-cocycle \(K\) (as in Theorem \ref{thm:xi-mixed-curvature-zero-criterion}).
Then there exists a natural integer \(2\)-cocycle \(\Omega\) (obtained by composing \(\alpha\) and \(K\) via the total differential in the corresponding bicomplex) satisfying:
\[
[\Omega]\in H^2(\mathcal B;\mathbb Z)
\]
is invariant under gauge transformations; and the pairing with any closed \(2\)-chain \(\Sigma\)
\[
\langle[\Omega],\Sigma\rangle\in\mathbb Z
\]
characterizes the net "time defect" over \(\Sigma\).
Therefore, all obstructions to time unification and additive-multiplicative unification can be compressed into a single integrable integer characteristic class.
\end{theorem}

\begin{proof}
Both \(\alpha\) and \(K\) are integer-valued \(2\)-cocycles; their gauge transformations introduce only integer coboundary terms, so the cohomology class \([\Omega]\) of the total cocycle induced in the total complex is an invariant.
The pairing with closed \(2\)-chains follows from the fundamental properties of cohomology and takes integer values; \([\Omega]=0\) if and only if both the composition anomaly and mixed curvature can be simultaneously eliminated, thereby yielding unified time and unified composition structure.
\end{proof}

\begin{theorem}[Time Index Theorem: The \(K_1\) Pairing of the Time Spectral Triple Equals the Time Holonomy]\label{thm:xi-time-index-theorem-k1-holonomy}
Under the setting of Theorem \ref{thm:xi-time-spectral-triple-distance-volume}, let \((\mathcal A,\mathcal H,D)\) be the spectral triple constructed from the time cocycle \(\widetilde\ell:\mathcal G\to\mathbb Z\) (\(D\delta_g=\widetilde\ell(g)\delta_g\)), and let \(P:=\mathbf 1_{[0,\infty)}(D)\) be the nonnegative spectral projection of \(D\).
For any unitary \(u\in\mathcal A\) (representing a closed loop class in \(K_1(\mathcal A)\)), the Fredholm index satisfies
\[
\mathrm{Index}(PuP)=\langle [u],[\tau_{\widetilde\ell}]\rangle,
\]
where \([\tau_{\widetilde\ell}]\) is the cohomology class of the cyclic \(1\)-cocycle induced by the time cocycle.
Moreover, in the groupoid closed-loop formulation, this pairing degenerates to the net time holonomy along the closed loop; equivalently, it agrees with the pairing of the time characteristic class \([\Omega]\) from Theorem \ref{thm:xi-time-characteristic-class-integrality} on the corresponding \(2\)-chain.
\end{theorem}

\begin{proof}
The index equality is the standard Connes pairing formula for odd spectral triples: \(\mathrm{Index}(PuP)\) is given by the pairing of the Chern character in \(K_1(\mathcal A)\) with cyclic cohomology.
In the present construction, \(D\) is diagonalized on the groupoid basis, so the corresponding cyclic cocycle can be taken as the groupoid cohomological representative of the time cocycle; evaluation on a closed loop gives the net increment of the time cocycle along the loop, which is thereby equivalent to the holonomy.
The consistency with \([\Omega]\) follows from the fact that when the "closed loop boundary is zero," the net increment can be rewritten as the integral pairing of the two-dimensional characteristic class on a filling \(2\)-chain.
\end{proof}

\paragraph{Time Geometry and Auditable Complexity}

\begin{theorem}[Time Homotopy Invariant for the Audit Equality Problem: The Certificate Dehn Function Is Uniquely Determined by the Time Geometry]\label{thm:xi-audit-dehn-function-time-geometry}
Suppose the protocol is given by a finitely generated convergent rewriting system, and let \(\delta(n)\) be the Dehn-type function: for any two words \(w\sim w'\) of length at most \(n\), \(\delta(n)\) is the supremum of the worst-case time length of the shortest audit certificate.
If the elementary step time is strictly additive and homotopy-invariant with respect to \(2\)-certificates, then \(\delta(n)\) is quasi-isometrically equivalent to the isoperimetric function of the time-metric Cayley graph.
In particular,
\[
\delta(n)=O(n)
\quad\Longleftrightarrow\quad
\text{the time-metric Cayley graph satisfies a linear isoperimetric inequality (hence is Gromov hyperbolic)},
\]
and consequently certificate verification admits a linear audit bound.
\end{theorem}

\begin{proof}
The definition of \(\delta(n)\) precisely externalizes the "minimal certificate area/time of word equivalence" as an isoperimetric function; strict additivity and homotopy invariance of time ensure that this function depends only on the quasi-isometry type of the time metric.
The equivalence between linear isoperimetric inequality and Gromov hyperbolicity is a standard result in geometric group theory; in the present setting, \(\delta(n)=O(n)\) yields a linear isoperimetric bound, which implies hyperbolicity and a linear audit bound.
\end{proof}

\paragraph{Growth Threshold, Maximal Visible Capacity, and Disambiguation Time Decomposition}

\begin{theorem}[Time Phase Transition Threshold Determined by the Growth Exponent, with Central Defect Rank Controlling the Dimension of the Equilibrium State Family]\label{thm:xi-time-phase-transition-growth-and-defect-rank}
On an ergodic transitive component, take the forward composable semigroup (or positive cone) generated by elementary steps
\[
\mathcal G_+:=\{g:\widetilde{\ell}(g)\in\mathbb N\},
\]
and define the time ball counting exponent
\[
\kappa:=\lim_{t\to\infty}\frac{1}{t}\log \#\{g\in\mathcal G_+:\widetilde{\ell}(g)\le t\},
\]
and define the time partition function (with source fixed)
\[
Z(\lambda):=\sum_{g\in\mathcal G_+} e^{-\lambda \widetilde{\ell}(g)}.
\]
Then there exists a critical value \(\lambda_c=\kappa\) such that \(Z(\lambda)<\infty\) if and only if \(\lambda>\lambda_c\).
Moreover, in the region \(\lambda>\lambda_c\), the equilibrium state compatible with this weight is unique; at the critical point \(\lambda=\lambda_c\), the dimension of the extremal decomposition of equilibrium states is controlled by the rank of the scale--time holonomy defect group \(\mathfrak H\) (Theorem \ref{thm:xi-scale-time-connection-holonomy-criterion}): the affine dimension of the extremal ray family equals \(\mathrm{rank}(\mathfrak H)\).
\end{theorem}

\begin{proof}
By the exponential growth of the time ball count and the comparison test, the convergence threshold of \(\sum_{g\in\mathcal G_+} e^{-\lambda\widetilde\ell(g)}\) is determined by the exponent \(\kappa\), hence \(\lambda_c=\kappa\).
The uniqueness and structural source of the critical degeneracy stem from the bifurcation of "whether the anomaly/defect can be absorbed into the center": when \(\lambda>\lambda_c\), the Gibbs weight yields a unique equilibrium state; when \(\lambda=\lambda_c\), nontrivial holonomy produces central parameter freedoms in the equivalence classes, so that equilibrium states form an affine family parametrized by the Pontryagin dual of \(\mathfrak H\), whose dimension is given by \(\mathrm{rank}(\mathfrak H)\).
\end{proof}

\begin{theorem}[Maximality Principle of Zeckendorf's Forbidden \(11\): The Maximal Visible Capacity among Locally Invertible Normalization Classes Is Uniquely Realized by the Golden Shift]\label{thm:xi-zeckendorf-golden-shift-max-capacity}
Consider all binary shift constraints \(Y\subset\{0,1\}^{\mathbb Z}\) satisfying: there exists a finite-radius, invertible, and confluent local normalization such that every finite word can be locally normalized to a unique normal form, and the time length defined by the normal form is strictly additive under composition.
Then among such constraints, the topological entropy \(h_{\mathrm{top}}(Y)\) attains a maximum, uniquely (up to conjugacy) realized by the golden shift: its adjacency matrix is
\[
\begin{pmatrix}
1 & 1\\[2pt]
1 & 0
\end{pmatrix},
\]
equivalent to the unique forbidden word being \(11\).
Therefore, the Zeckendorf grammar can be derived from the three principles of "time additivity + local invertibility + maximal visible capacity," rather than being assumed a priori.
\end{theorem}

\begin{proof}
Local invertibility and confluence yield uniqueness of the normal form language; strict additivity of time promotes the normal form length to a composable cocycle.
Under the binary local normalization framework, the counting growth of legal grammars must be compatible with the Fibonacci recurrence (cf.\ the weight rigidity framework of Proposition \ref{prop:xi-local-binary-normalization-fibonacci-rigidity}), so that the entropy rate is upper-bounded by the Perron--Frobenius root \(\varphi\), with equality if and only if the grammar degenerates to the golden mean constraint (forbidden \(11\)).
Conjugacy uniqueness follows from the standard classification of binary irreducible SFTs at a given PF root.
\end{proof}

\begin{proposition}[Exact Decomposition of the Minimum Disambiguation Time: The Geometric Term Is \(\mathbb E\lbrack \log\mathrm{Ind}\rbrack\), and the Statistical Term Is the Intra-Fiber Relative Entropy]\label{prop:xi-min-disambiguation-time-entropy-decomposition}
Let the fold \(\mathrm{Fold}:\Omega\to X\) induce fiber index \(\mathrm{Ind}(x):=|\mathrm{Fold}^{-1}(x)|\), with a given conditional distribution \(\mu(\cdot\mid x)\).
Denote the intra-fiber uniform distribution by \(u_x\).
Then the conditional entropy decomposition identity is
\[
H_\mu(\omega\mid x)
=\mathbb E_\mu[\log \mathrm{Ind}(x)]
-\mathbb E_\mu\!\bigl[D_{\mathrm{KL}}(\mu(\cdot\mid x)\,\|\,u_x)\bigr].
\]
Furthermore, for any sequential readout protocol with alphabet size \(|\mathcal A|\), the minimum expected disambiguation time satisfies
\[
\inf \mathbb E_\mu[T]
=
\frac{H_\mu(\omega\mid x)}{\log|\mathcal A|}+O(1)
=
\frac{\mathbb E_\mu[\log \mathrm{Ind}(x)]}{\log|\mathcal A|}
-\frac{\mathbb E_\mu[D_{\mathrm{KL}}(\mu(\cdot\mid x)\,\|\,u_x)]}{\log|\mathcal A|}
+O(1).
\]
\end{proposition}

\begin{proof}
For each \(x\),
\(
D_{\mathrm{KL}}(\mu(\cdot\mid x)\,\|\,u_x)=\log\mathrm{Ind}(x)-H(\mu(\cdot\mid x))
\);
taking the expectation over \(x\) yields the first identity.
The second identity follows directly from the optimal stopping time bound of Theorem \ref{thm:xi-optimal-recovery-time-equals-conditional-entropy}.
\end{proof}

\paragraph{Time Order Structure and Realization Rigidity}

\begin{proposition}[Time Positive Ordering and Torsion-Freeness: A Strictly Positive Time Cocycle Forces Left-Orderability and Absence of Finite-Order Elements]\label{prop:xi-time-positive-cocycle-left-order-torsionfree}
Suppose that on the isotropy groups (homotopy groups) of each fixed point in a transitive component, the time cocycle \(\widetilde{\ell}\) satisfies
\(\widetilde{\ell}(gh)=\widetilde{\ell}(g)+\widetilde{\ell}(h)\),
and for any nontrivial element \(g\neq e\), \(\widetilde{\ell}(g)\neq 0\).
Define the positive cone
\[
P:=\{g:\widetilde{\ell}(g)>0\}.
\]
Then \(P\) is a subsemigroup disjoint from its inverse, so that the isotropy group is left-orderable; moreover, no torsion elements exist (if \(g^n=e\) then \(g=e\)).
\end{proposition}

\begin{proof}
By additivity, \(P\cdot P\subseteq P\), and \(g\in P\Rightarrow g^{-1}\notin P\) (since \(\widetilde{\ell}(g^{-1})=-\widetilde{\ell}(g)\)).
From this, one defines a left-invariant total order: \(g\prec h\) if and only if \(h g^{-1}\in P\).
If \(g^n=e\), then \(0=\widetilde{\ell}(g^n)=n\widetilde{\ell}(g)\), and by nondegeneracy \(\widetilde{\ell}(g)=0\), hence \(g=e\).
\end{proof}

\begin{theorem}[Time Conjugacy Rigidity: The Same Time Cocycle Determines a Bounded Conjugacy Class of Finite-State Realizations]\label{thm:xi-time-conjugacy-rigidity-bounded-class}
Suppose two finite-state realizations \((\Gamma,\widetilde{\ell})\), \((\Gamma',\widetilde{\ell}')\) represent the same protocol (two coverings of the same visible groupoid), and satisfy time cocycle cohomological equivalence: there exists a vertex function \(b\) such that
\[
\widetilde{\ell}'=\widetilde{\ell}+\partial b.
\]
Then there exists a sliding block conjugacy \(\Phi\) (with finite memory and finite anticipation) sending the path space of \(\Gamma\) to that of \(\Gamma'\), satisfying strict time preservation
\[
\widetilde{\ell}'(\Phi(\gamma))=\widetilde{\ell}(\gamma)
\quad\text{for almost every path \(\gamma\)}.
\]
Furthermore, the conjugacy window size admits a uniform upper bound given by the maximum absolute value of the scale--time holonomy generators (viewed as the "time curvature norm").
\end{theorem}

\begin{proof}
Cocycle cohomological equivalence means the time weights of the two realizations differ only by a coboundary of a vertex potential; for the cumulative time along a path, this coboundary telescopes, so that under appropriate state relabeling/state splitting recoding, a strictly time-preserving conjugacy can be realized.
The boundedness of the window size comes from the finite amplitude control of \(b\), and the necessary amplitude of \(b\) is lower-bounded by the maximum absolute value of the holonomy (curvature) generators.
\end{proof}

\begin{proposition}[Identifiability of the Time Partition Function: The Fiber Index Multiset at Each Level Can Be Reconstructed from \(\{Z_m(\lambda)\}\) Alone]\label{prop:xi-time-partition-identifiability-ind-multiset}
At each scale \(m\), let the fiber index multiset be \(\{\mathrm{Ind}_m(x):x\in X_m\}\subset\mathbb N\), and define the time potential \(T_m(x)=\log \mathrm{Ind}_m(x)\).
Consider the Mellin--Laplace-type partition function
\[
Z_m(\lambda):=\sum_{x\in X_m} e^{-\lambda T_m(x)}=\sum_{x\in X_m}\mathrm{Ind}_m(x)^{-\lambda}.
\]
If at a fixed scale \(m\), \(\mathrm{Ind}_m(x)\) takes only finitely many integer values, then the values of \(Z_m(\lambda)\) on any set of \(\lambda\) containing an accumulation point uniquely determine this finite discrete measure, and hence uniquely determine the multiplicity of each index value.
Furthermore, when \(m\to\infty\) and the pressure limit
\(
P(\lambda)=\lim_{m\to\infty}\frac{1}{m}\log Z_m(\lambda)
\)
exists and is differentiable on an interval, the Legendre dual of \(P\) uniquely gives the multifractal spectrum of fiber time (cf.\ Theorem \ref{thm:xi-fiber-time-multifractal-legendre}).
\end{proposition}

\begin{proof}
Under the finite spectrum assumption, \(Z_m(\lambda)\) is a linear combination of finitely many functions \(\lambda\mapsto n^{-\lambda}\); these functions are linearly independent in the real-analytic sense, so equality on any set containing an accumulation point implies coefficient-by-coefficient (i.e., multiplicity-by-multiplicity) equality.
The second part follows directly from the large deviation--Legendre duality mechanism given by the pressure limit and the G\"artner--Ellis condition.
\end{proof}
